%Binaire, SL 11, Faynot


\setcounter{equation}{0}
\setcounter{Q}{0}

\begin{center}
\begin{minipage}{0.48\textwidth}
\includegraphics[width=\textwidth]{/home/jeff/Documents/IUTOrsay/1-Cours/4p13-Alliages/Cours-TD/Exos/Binaires-Lecture/Heterogene-Ex01/figures/Fig01}
\end{minipage}\hfill\begin{minipage}{0.48\textwidth}
\captionof{figure}{Diagramme binaire isobare, à une pression de \SI{1}{\bar}, du mélange \ce{Au}-\ce{Si}.}
[La composition est donnée par la fraction massique en silicium dans le mélange]
\end{minipage}
\end{center}

Le diagramme binaire isobare ($P = 1 ~ \si{\bar}$) de l’or \ce{Au} et du silicium \ce{Si} est donné ci-contre, la composition étant exprimée en fraction molaire de \ce{Si}, soit $x_{\ce{Si}}$.

\Q Identifier les courbes du solidus et du liquidus.

\Q Calculer la variance pour les différents
systèmes représentés sur le diagramme.

\Q Un mélange liquide d’or et de silicium, de composition $x_{\ce{Si}} = 0.70$ est refroidi. Déterminer la température de solidification commençante et la nature du solide qui apparaît, puis tracer la courbe de refroidissement correspondante, entre les températures \SI{1700}{\celsius} et \SI{300}{\celsius}.

\Q Un mélange équimolaire d’or et de silicium, de masse $m = 200 ~ \si{\gram}$, est porté à \SI{700}{\celsius}.

\begin{enumerate}
\item Déterminer la nature, la composition en fractions molaires et les quantités de matière présentes dans les phases.

\item De l’or est progressivement ajouté à ce mélange, tout en maintenant la température à \SI{700}{\celsius}. Déterminer la quantité (\si{\mole}) d’or à ajouter pour que le système devienne monophasé.

\end{enumerate}

\paragraph*{Données} 
\begin{align*}
M(\ce{Si}) &= 28 ~ \si{\gram\per\mole} ; \\
M(\ce{Au}) &= 197 ~ \si{\gram\per\mole} .
\end{align*}
